% Section 7.3: Kriptanalisis Knapsack
\section{Kriptanalisis dan Kelemahan Sistem Knapsack}

\vspace{0.5cm}
\noindent\colorbox{blue!10}{\parbox{\dimexpr\textwidth-2\fboxsep}{\textbf{\large Sub-CPMK yang Dicakup dalam Bab Ini:}}}
\vspace{0.3cm}
\begin{itemize}[leftmargin=*, itemsep=5pt]
Mahasiswa mampu menganalisis kelemahan sistem Merkle-Hellman melalui serangan Adi Shamir dan menarik kesimpulan mengenai keamanan sistem tersebut.
\end{itemize}
\vspace{0.5cm}

Meskipun Subset Sum Problem secara umum adalah NP-complete, sistem Merkle-Hellman memiliki struktur yang dapat dieksploitasi. Pada tahun 1982, Adi Shamir menunjukkan bahwa sistem ini tidak aman.

\subsection{Serangan Adi Shamir}
Adi Shamir menggunakan teknik pemrograman linear dan analisis lattice (kisi) untuk memecahkan sistem ini. Inti dari serangannya adalah menemukan nilai $m$ dan $q$ (atau nilai yang setara) yang dapat mengubah barisan publik $P$ kembali menjadi barisan superincreasing.

Shamir membuktikan bahwa pemilihan $m$ dan $q$ tidak cukup untuk menyembunyikan sifat superincreasing dari barisan $S$. Dengan menganalisis hubungan antara elemen-elemen dalam kunci publik $P$, penyerang dapat mengestimasi nilai $m$ dan $q$ dan kemudian melakukan dekripsi tanpa memiliki kunci privat asli.

\subsection{Kondisi Keamanan Saat Ini}
Akibat serangan Shamir dan pengembangan algoritma LLL (\textit{Lenstra-Lenstra-Lovász}), hampir semua varian kriptografi Knapsack yang menggunakan barisan superincreasing telah dianggap tidak aman. Saat ini, algoritma ini tidak digunakan dalam sistem produksi, namun tetap dipelajari sebagai contoh penting dalam sejarah kriptografi mengenai bagaimana sebuah masalah yang "sulit secara umum" (NP-complete) tidak selalu menjamin keamanan jika implementasinya memberikan struktur yang bisa dieksploitasi.

\subsection{Perbandingan Algoritma Asimetris}
Sebagai penutup bab ini, berikut adalah ringkasan perbandingan antara sistem asimetris utama yang telah kita pelajari.

\begin{table}[h]
\centering
\begin{tabular}{|l|l|l|l|}
\hline
\textbf{Algoritma} & \textbf{Dasar Keamanan} & \textbf{Kecepatan} & \textbf{Status} \\ \hline
RSA & Faktorisasi Prima & Sedang & Aman (Kunci Besar) \\ \hline
DH / ElGamal & Logaritma Diskrit & Sedang & Aman (Kunci Besar) \\ \hline
ECC & ECDLP & Sangat Cepat & Sangat Aman \\ \hline
Knapsack & Subset Sum & Cepat & \textbf{Tidak Aman} \\ \hline
\end{tabular}
\caption{Perbandingan Algoritma Kunci Publik}
\end{table}

\begin{aktivitas}
Baca artikel mengenai Algoritma LLL dan diskusikan bagaimana teknik analisis lattice dapat memecahkan masalah Knapsack secara umum.
\end{aktivitas}

\begin{latihan}
Apa perbedaan mendasar antara keamanan RSA dan keamanan Knapsack Merkle-Hellman? Mengapa satu tetap bertahan sementara yang lain runtuh?
\end{latihan}

